% !Mode:: "TeX:UTF-8"

%%% 说明: 此部分需要自己填写的内容:  (1) 攻博(硕)期间科研成果 (2) 致谢

%%%%%%%%%%%%%%%%%%%%%%%
%%%------- 攻博(硕)期间科研成果 -------%%%
%%%%%%%%%%%%%%%%%%%%%%%
\reseachresult

\zihao{5}
\begin{enumerate}[{[1]}]
\item   Deep Ritz method for elliptical multiple eigenvalue problems. \textit{Journal of Scientific Computing}, 2024, 98(2), 48.

\item   Analysis of deep Ritz methods for semilinear elliptic equations. \textit{Numerical Mathematics Theory Methods and Applications}, 2024, 17(1), 181-209. 

\item   Potential identification via Tikhonov-PINNs. \textit{Inverse Problems}, 2025, 41(11): 115008.

\end{enumerate}


%%%%%%%%%%%%%%%%%%%%%%%
%%% --------------- 致谢 ------------- - %%%
%%%%%%%%%%%%%%%%%%%%%%%
\acknowledgement



岁月匆匆，五年的博士求学生涯如白驹过隙.回首来路，这段时光虽也有过迷茫、困顿与无助的时刻，但整体的色调始终是积极、乐观、明亮的.无数次办公室里的灯火通明，下班路上的披星戴月，见证了我一路走来的轨迹——脚步虽慢，却走得踏实而坚定.

首先，最要感谢的是在五年时光里给予我悉心指导、领我踏上学术正途的几位恩师.衷心感谢我的导师吕老师，先生在同学们中有着极高的口碑，从论文的选题立意到最终的修改定稿，吕老师都倾注了大量心血。老师严谨的治学精神和务实的作风，深刻地影响了我；感谢杨志坚老师，杨老师严谨的治学态度、渊博的专业知识以及对学术前沿的敏锐洞察力，让我在求学路上受益匪浅；还要感谢焦雨领老师，焦老师在许多具体的学术问题上给予了我极其细致的指导，极大地拓宽了我的学术视野，提升了我的学术品位.

其次，感谢五年里朝夕相伴的同门师兄弟们.在导师们积极正面的价值观引领下，课题组始终保持着团结互助、齐心协力的良好氛围.在这里，我感受到的是携手并进的勇气，而非相互竞争的内耗.特别感谢王封儒、段晨光等师兄，是你们的帮助让我这名“学术萌新”顺利步入正轨；感谢袁成等师兄，你们的鼓励让我坚定了求学问道的决心；感谢715办公室的同门们，朝夕的陪伴是忙碌科研生活中的温馨调剂，你们的勤奋努力也时常化作我奋发向前的动力.

此外，还要深深感谢在五年中陪伴我、给予我无尽力量的家人们.最要感谢的是我的妻子，我们相识于校园，她携手陪我走过了整个博士生涯.她的存在，是我身处低谷时依然心怀微光的理由，也是我顺境中牢记责任的源泉.同样要感谢我的父母，家庭的无私托举，是我能够抵御外在风雨、安心潜行于学术道路的坚实后盾.

一路走来，我深知今日的欢欣绝非我一人之功.短短一页纸，难以写尽所有应该感谢的名字；寥寥数笔，亦无法穷尽心中的万语千言.唯有将师长亲友的激励化作前进的动力，将大家的期望挑作肩头的责任.前路虽或曲折漫长，但我坚信，只要笃定前行，终将拨云见日.行笔至此，唯以东坡先生的半阕《定风波》自勉：

“莫听穿林打叶声，何妨吟啸且徐行.竹杖芒鞋轻胜马，谁怕？一蓑烟雨任平生.”











%%%%%%%%%%%%%%%%%%%%%%%%%%%%%%%%%%%%%%%
%%%%%%%--判断是否需要空白页-----------------------------
  \iflib
  \else
  \newpage
  %\thispagestyle{empty}
  \cleardoublepage
  \fi
%%%%%%%-------------------------------------------------

